\documentclass[11pt]{article}

\usepackage[a4paper,margin=1in]{geometry}
\usepackage{amsmath,amssymb,amsthm,mathtools}
\usepackage{booktabs}
\usepackage{array}
\usepackage{microtype}
\usepackage{xcolor}
\usepackage{hyperref}
\usepackage{enumitem}

\hypersetup{
  colorlinks=true,
  linkcolor=blue!55!black,
  citecolor=blue!55!black,
  urlcolor=blue!55!black
}

\newtheorem{definition}{Definition}
\newtheorem{proposition}{Proposition}
\newtheorem{theorem}{Theorem}
\newtheorem{remark}{Remark}

\newcommand{\ii}{\mathrm{i}}
\newcommand{\dd}{\mathrm{d}}
\newcommand{\E}{\mathbb{E}}
\newcommand{\R}{\mathbb{R}}
\newcommand{\C}{\mathbb{C}}
\newcommand{\Tr}{\operatorname{Tr}}
\newcommand{\rank}{\operatorname{rank}}
\newcommand{\Span}{\operatorname{span}}
\newcommand{\cF}{\mathcal{F}}
\newcommand{\cI}{\mathcal{I}}
\newcommand{\cG}{\mathcal{G}}
\newcommand{\cC}{\mathcal{C}}
\newcommand{\cR}{\mathcal{R}}

\title{\textbf{Covariant Local Jets of Response-Matched Neutral-Atom Controls:\\
From Quartic Prediction to Formal Finite-Error Certification}}

\author{Y.\,Y.\,N.\ Li\\
\small Independent Researcher}

\date{\today}

\begin{document}
\maketitle

\begin{abstract}
Control sequences can agree at their ideal endpoint and in their complete
first-order response to specified implementation errors while differing at
finite error.  We study this phenomenon in a two-atom Rydberg model with
piecewise-constant phase control.  Twenty-four control phases are constrained
to match the output state and its first-order state response to global
amplitude, detuning, and interaction errors.  The resulting numerical
constraint Jacobian has rank sixteen, leaving an eight-dimensional
implementation fibre.  Paths on this fibre display distinct finite-error
infidelities.

We first identify a fourth-order symmetric response coefficient,
$\cG_4$, as a strong low-order predictor of mean finite-error performance.
In a protocol-frozen prospective twenty-path validation it attains Spearman correlation
$0.998496$ and predicts the best path correctly.  We then recover the full
three-dimensional symmetric fourth-order tensor.  Under a nonorthogonal
change of error coordinates, direct recomputation agrees with the tensor
transformation law to $1.1\times10^{-14}$, while contraction with the
transformed fourth noise moment is invariant to $2.0\times10^{-15}$.

The quartic term is not universally sufficient.  We therefore compute the
zero-error Taylor jet through order thirty by block-matrix exponentials and
bound the uncomputed tail analytically.  Finally, a 192-bit Arb
midpoint-radius calculation uses outward-rounded complex balls, a
64-point Cauchy extraction, a rigorous alias enclosure, and an analytic
order-32 tail enclosure.  For a frozen twelve-path cohort, the quartic term
alone certifies 35 of 66 path pairs, whereas the order-thirty jet certifies
all 66 pairs.  All ordinary held-out outcomes lie inside the formal
intervals.

We then close the numerical-matching gap.  A square sixteen-dimensional
transverse chart and outward-rounded Krawczyk calculation certify one locally
unique exact state-and-first-response-matched root inside each of twelve
declared phase boxes.  Direct 192-bit Arb propagation of those boxes through
the six finite-error points certifies the frozen order of all 66 pairs.  The
order-thirty jet propagated over the same boxes separately certifies 42 pairs,
all in the frozen direction, with the other 24 unresolved through interval
widening.

The result is a formal exact-root finite-error ordering certificate for a
serialized-decimal, finite-dimensional Hamiltonian model.  It is not evidence
from PASQAL hardware and does not establish global uniqueness or a global
eight-dimensional fibre theorem.
\end{abstract}

\section{Introduction}

Pulse-level control of programmable neutral-atom arrays provides direct
access to Rabi amplitude, phase, detuning, geometry, and Rydberg-mediated
interactions \cite{Silverio2022,Henriet2020,Saffman2016}.  This flexibility
also creates a basic identification problem.  Two implementations may
realize the same ideal state and have the same leading response to calibrated
errors, yet behave differently when the errors are finite.  Endpoint
agreement alone does not identify an implementation, and even first-order
robustness data may leave a nontrivial family of controls.

Composite-pulse and robust-control methods commonly cancel or suppress
leading error terms \cite{Wimperis1994,Kukita2022,Barnes2015}.  The question
asked here is different:
\begin{quote}
\emph{After the ideal output state and the complete first-order output-state
response have been matched, what local object predicts the residual
finite-error ordering inside the remaining implementation fibre?}
\end{quote}

We address the question in four stages.
\begin{enumerate}[leftmargin=2.2em]
\item We construct phase schedules with a common endpoint and common
first-order state response to amplitude, detuning, and Rydberg-interaction
errors.
\item We test a low-order quartic response scalar, $\cG_4$, using a frozen
ranking followed by held-out finite-error evaluation.
\item We promote $\cG_4$ to a coordinate-covariant fourth-order tensor whose
contraction with a physical noise moment is invariant.
\item We compute higher local derivatives and produce an outward-rounded
Arb ball-arithmetic certificate of finite-error ordering.
\item We certify exact matched representatives by Krawczyk inclusions and
propagate their full phase boxes directly to the finite-error functional.
\end{enumerate}

The central outcome is deliberately narrower than a universal robustness
law.  The fourth-order term is often predictive and can certify a substantial
subset of pairwise comparisons, but close paths require higher derivatives.
For the serialized model studied here, the order-thirty local jet plus an
analytic tail bound freezes a complete finite-error ordering.  Krawczyk
inclusions and direct propagation then certify that the same order holds at
the exact matched roots.  This distinguishes statements that are often
conflated:
\[
\begin{aligned}
\text{low-order correlation}
&\ne \text{partial pairwise certification}\\
&\ne \text{complete bounded reconstruction}\\
&\ne \text{exact-root finite-radius certification}.
\end{aligned}
\]

\section{Two-atom control model}

\subsection{Hamiltonian}

We use the ordered basis
\[
\{|gg\rangle,|gr\rangle,|rg\rangle,|rr\rangle\}.
\]
Let
\[
X_\Sigma=X\otimes I+I\otimes X,\qquad
Y_\Sigma=Y\otimes I+I\otimes Y,
\]
\[
N_\Sigma=n\otimes I+I\otimes n,\qquad
N_{rr}=n\otimes n,
\]
where $n=|r\rangle\langle r|$.  During segment $j$, the Hamiltonian is
\begin{align}
H_j(\boldsymbol{\epsilon})
={}&
\frac{\Omega}{2}(1+\epsilon_\Omega)
\left(\cos\phi_j\,X_\Sigma+\sin\phi_j\,Y_\Sigma\right)
\nonumber\\
&-\Omega\,\epsilon_\Delta N_\Sigma
+V(1+\epsilon_V)N_{rr}.
\label{eq:Hamiltonian}
\end{align}
The dimensionless error vector is
\[
\boldsymbol{\epsilon}
=(\epsilon_\Omega,\epsilon_\Delta,\epsilon_V).
\]
The nominal parameters are
\[
\Omega=2\pi~\mathrm{rad}/\mu\mathrm{s},\qquad
V=4\Omega,
\]
with 24 segments of duration
\[
\tau=0.1~\mu\mathrm{s},
\qquad
T=24\tau=2.4~\mu\mathrm{s}.
\]
The nominal interaction can be represented as $V=C_6/R^6$; using the
Pulser DigitalAnalogDevice coefficient in the accompanying implementation
gives a nominal separation of approximately $7.744~\mu\mathrm{m}$.

For a control path
\[
\gamma=(\phi_1,\ldots,\phi_{24})\in\mathbb{T}^{24},
\]
the output state is
\begin{equation}
|\psi_\gamma(\boldsymbol{\epsilon})\rangle
=
\prod_{j=24}^{1}
\exp[-\ii\tau H_j(\boldsymbol{\epsilon})]
|gg\rangle.
\label{eq:state}
\end{equation}
All primary calculations use exact matrix exponentials for the
piecewise constant, finite-dimensional model.  Pulser/QutipEmulator was used
in an earlier two-atom cross-check, but the formal certificate in this work
is independent of Pulser and does not invoke a cloud backend.

\subsection{Loss and held-out error law}

Let $\gamma_{\mathrm{ref}}$ denote the reference schedule and define
\[
|\psi_{\mathrm{ref}}\rangle
=|\psi_{\gamma_{\mathrm{ref}}}(\mathbf{0})\rangle.
\]
The state infidelity is
\begin{equation}
\cI_\gamma(\boldsymbol{\epsilon})
=1-\left|
\langle\psi_{\mathrm{ref}}|
\psi_\gamma(\boldsymbol{\epsilon})\rangle
\right|^2.
\label{eq:infidelity}
\end{equation}

The held-out error design contains the six signed axial errors
\begin{equation}
\mathcal{E}
=
\left\{
\pm0.06\,e_\Omega,\,
\pm0.04\,e_\Delta,\,
\pm0.05\,e_V
\right\}.
\label{eq:errorlaw}
\end{equation}
The primary performance functional is the mean
\begin{equation}
\overline{\cI}_\gamma
=\frac{1}{6}
\sum_{\boldsymbol{\epsilon}\in\mathcal{E}}
\cI_\gamma(\boldsymbol{\epsilon}).
\label{eq:meanloss}
\end{equation}
Worst-case infidelity is retained as a secondary diagnostic.  The present
predictors do not reliably rank the worst case, and no worst-case theorem is
claimed.

\section{Matched implementation fibre}

\subsection{Endpoint and first-order constraints}

At zero error, remove the unobservable global phase by the horizontal
projector
\[
P_\perp=I-|\psi_0\rangle\langle\psi_0|.
\]
Define the horizontal first-order responses
\begin{equation}
J_{\gamma,a}
=
P_\perp
\left.
\frac{\partial}{\partial\epsilon_a}
|\psi_\gamma(\boldsymbol{\epsilon})\rangle
\right|_{\boldsymbol{\epsilon}=0},
\qquad
a\in\{\Omega,\Delta,V\}.
\label{eq:tangents}
\end{equation}

\begin{definition}[Numerically matched implementation fibre]
For a declared tolerance $\eta$, the numerical matched fibre
$\cF_\eta$ consists of controls $\gamma$ satisfying
\begin{align}
\cI_\gamma(\mathbf{0})&\le\eta,
\nonumber\\
\|J_{\gamma,a}-J_{\mathrm{ref},a}\|&\le\eta,
\qquad a\in\{\Omega,\Delta,V\}.
\label{eq:fibre}
\end{align}
\end{definition}

The actual optimizer retains the full phase-aligned complex output state and
the three full phase-aligned complex state tangents.  Redundant real
components are retained in the residual vector, while singular-value
decomposition determines the independent rank.

\subsection{Dimension count and generation}

The control has 24 phase variables.  At the reference solution, the numerical
constraint Jacobian has rank
\[
\rank D\cC(\gamma_{\mathrm{ref}})=16,
\]
leaving an eight-dimensional numerical null space.  Candidate paths are
generated by:
\begin{enumerate}[leftmargin=2.2em]
\item drawing a seeded direction in the numerical null space;
\item applying a finite displacement;
\item projecting back to the constraint set by nonlinear least squares;
\item rejecting candidates whose constraint residual or pairwise phase
distance fails a predeclared gate.
\end{enumerate}
No held-out finite-error outcome is evaluated during candidate generation.
The maximum constraint residual in the reported runs is of order
$10^{-11}$.

\begin{theorem}[Certified local exact representatives]
\label{thm:krawczyk}
Fix the eight fibre coordinates used to define each member of the declared
twelve-path cohort and use the frozen sixteen-dimensional transverse chart.
For every path, the outward-rounded Krawczyk image lies strictly inside one
of the declared transverse boxes.  Consequently, each box contains one
locally unique zero of the exact state-and-first-response constraint map.
\end{theorem}

\begin{proof}[Computer-assisted proof]
The constraint dimension is sixteen.  Exact Newton refinement precedes the
formal audit, after which the phase/error Jacobian is evaluated with
192-bit Arb balls.  A frozen radius schedule is tested without reference to
finite-error outcomes.  Strict Krawczyk inclusion passes for all twelve
paths, at radii $3\times10^{-12}$ for paths \texttt{pv01}--\texttt{pv10}
and $10^{-11}$ for \texttt{pv11}--\texttt{pv12}.
\end{proof}

\begin{remark}[Local, not global, fibre statement]
Theorem~\ref{thm:krawczyk} proves existence and uniqueness in each declared
transverse box at fixed fibre coordinates.  It does not by itself establish
global uniqueness, connectedness, or a global eight-dimensional manifold
structure for the complete matched set.
\end{remark}

\section{Local response hierarchy}

\subsection{Symmetric local response}

Introduce normalized error coordinates $z\in\R^3$ through
\[
\boldsymbol{\epsilon}=Dz,\qquad
D=\operatorname{diag}(0.06,0.04,0.05).
\]
For a direction $v\in\R^3$, define the symmetric response
\begin{equation}
S_\gamma(t,v)
=
\frac{
\cI_\gamma(tDv)+
\cI_\gamma(-tDv)
}{2}.
\label{eq:symmetric}
\end{equation}
Analyticity of the finite matrix exponential gives
\begin{equation}
S_\gamma(t,v)
=
q_{2,\gamma}(v)t^2
+q_{4,\gamma}(v)t^4
+q_{6,\gamma}(v)t^6+\cdots.
\label{eq:jet}
\end{equation}
Odd powers cancel by construction.

First-order output-state matching fixes the leading quadratic loss to
numerical precision.  Across the validation cohorts, the relative spread of
the averaged quadratic coefficient is approximately $10^{-7}$.  The first
useful discriminator is therefore the quartic coefficient.

\subsection{Quartic tensor}

There exists a symmetric fourth-order tensor $A^{(4)}_\gamma$ such that
\begin{equation}
q_{4,\gamma}(v)
=
A^{(4)}_{\gamma,ijkl}
v_i v_j v_k v_l.
\label{eq:A4}
\end{equation}
In three dimensions a symmetric fourth-order tensor has
\[
\binom{3+4-1}{4}=15
\]
independent components.  We recover them from 36 seeded directions.  The
directional Taylor coefficients are computed directly at $t=0$ through
order six by block-matrix exponentials, rather than by fitting nonzero-radius
samples.

For the six-axis law in Eq.~\eqref{eq:errorlaw}, define the fourth moment
\[
M^{(4)}_{ijkl}
=\E[z_i z_j z_k z_l].
\]
The scalar quartic predictor is
\begin{equation}
\cG_4(\gamma;M^{(4)})
=A^{(4)}_{\gamma,ijkl}M^{(4)}_{ijkl}.
\label{eq:G4}
\end{equation}
For the equally weighted signed axes,
\[
\cG_4
=\frac{1}{3}
\left(
A^{(4)}_{\Omega\Omega\Omega\Omega}
+A^{(4)}_{\Delta\Delta\Delta\Delta}
+A^{(4)}_{VVVV}
\right).
\]

\subsection{Coordinate covariance}

\begin{proposition}[Covariance and invariant contraction]
Let $z=Ry$ for an invertible real matrix $R$.  Then
\begin{align}
A^{(4)}_{y,abcd}
&=
A^{(4)}_{z,ijkl}
R_{ia}R_{jb}R_{kc}R_{ld},
\label{eq:Atransform}\\
M^{(4)}_{y,abcd}
&=
(R^{-1})_{ai}(R^{-1})_{bj}
(R^{-1})_{ck}(R^{-1})_{dl}
M^{(4)}_{z,ijkl},
\label{eq:Mtransform}
\end{align}
and
\begin{equation}
A^{(4)}_y:M^{(4)}_y
=A^{(4)}_z:M^{(4)}_z.
\label{eq:invariance}
\end{equation}
\end{proposition}

\begin{proof}
Equation~\eqref{eq:Atransform} follows by substituting $z=Ry$ into the
quartic form.  Equation~\eqref{eq:Mtransform} follows from
$y=R^{-1}z$.  Contracting all four indices cancels each $R$ with the
corresponding $R^{-1}$, proving Eq.~\eqref{eq:invariance}.
\end{proof}

This proposition is algebraic.  The numerical audit additionally recomputes
the tensor directly in the transformed coordinates.  For a nonorthogonal
map of condition number approximately $1.82$, the direct-refit relative
error is $1.08\times10^{-14}$ and the maximum invariant-contraction relative
error is $1.99\times10^{-15}$.

\section{Prospective quartic ranking}

\subsection{Discovery and validation separation}

The initial four named controls were used only for discovery and orientation.
A twelve-path cohort was then used to compare candidate formulas.  The
smallest sufficient rule on that training cohort was
\[
\text{smaller }\cG_4
\quad\Longrightarrow\quad
\text{smaller mean finite-error infidelity}.
\]
The addition of the sixth-order coefficient did not improve the training
rank correlation, so the quartic rule was frozen.

A new twenty-path cohort was generated with a different seed.  For each
candidate, a ranking certificate containing phases, local coefficients,
and predicted order was written and SHA-256 hashed before the held-out
outcome evaluator was enabled.  The predeclared primary gate was
\[
\rho_{\mathrm{Spearman}}\ge0.95,
\]
together with correct top-path identification.

\subsection{Quartic validation result}

The twenty-path Colab run gave
\begin{equation}
\rho_{\mathrm{Spearman}}
\left(
-\cG_4,-\overline{\cI}
\right)
=0.998496,
\qquad
p=3.71\times10^{-24}.
\label{eq:g4rho}
\end{equation}
The predicted and observed top path agreed.  Only one adjacent pair was
inverted in the complete ranking.  Two local coefficient-extraction windows
had rank correlation one, and the maximum local fit residual was
$1.15\times10^{-10}$.

Subsequent cohorts show that $\cG_4$ is not universally sufficient.  Depending
on the generated matched paths, prospective quartic Spearman correlations
ranged from approximately $0.85$ to $1.00$.  This variation motivated the
bounded higher-order analysis rather than post hoc replacement of the
quartic formula.

\section{High-order local reconstruction}

\subsection{Exact zero-point coefficients}

For one segment, write
\[
U_j(t)=\exp(A_j+tB_j),
\qquad
A_j=-\ii\tau H_j(0).
\]
The coefficient of $t^n$ can be obtained from the $(0,n)$ block of the
exponential of an upper-bidiagonal block matrix with $A_j$ on the diagonal
and $B_j$ on the first superdiagonal.  Segment series are multiplied and
truncated to order $N=30$.  This produces the zero-point state jet without
finite-radius differencing.

For the symmetric mean loss,
\begin{equation}
\overline{\cI}_\gamma(t)
=
C_{0,\gamma}
+C_{2,\gamma}t^2
+\cG_{4,\gamma}t^4
+\sum_{m=3}^{15}
\cG_{2m,\gamma}t^{2m}
+\cR_{32,\gamma}(t).
\label{eq:order30}
\end{equation}
The small $C_0$ term is retained because the numerical candidate endpoint is
not symbolically identical to the reference endpoint.

\subsection{Analytic tail bound}

Let $K_a$ be the integrated spectral norm of the Hamiltonian perturbation
along axis $a$:
\[
K_a=\int_0^T\|\partial_t H_a(t)\|_2\,\dd t.
\]
For the normalized held-out axes,
\[
(K_\Omega,K_\Delta,K_V)
\approx(0.90478,1.20637,3.01593).
\]
A Dyson-series norm estimate yields
\begin{equation}
|\cR_{32,\gamma}(1)|
\le
\frac{1}{3}
\sum_{a\in\{\Omega,\Delta,V\}}
\sum_{\substack{n\ge32\\n\ \mathrm{even}}}
\frac{(2K_a)^n}{n!}.
\label{eq:tail}
\end{equation}
The floating-point order-thirty audit evaluates the right-hand side as
approximately $1.2343\times10^{-11}$.  The maximum difference between the
order-thirty reconstruction and ordinary held-out propagation was below
$10^{-15}$ in the reported runs.

\subsection{Two notions of certification}

To test the quartic term alone, define the quartic centre
\[
\mu_{\gamma}^{(4)}
=\overline C_0+\overline C_2+\cG_{4,\gamma},
\]
and place the complete known correction
\[
(C_{0,\gamma}-\overline C_0)
+(C_{2,\gamma}-\overline C_2)
+\sum_{m=3}^{15}\cG_{2m,\gamma}
\]
inside the interval radius together with the tail bound.  Two paths are
quartic-certified only if these conservative intervals are disjoint.

For the order-thirty certificate, the centre includes all known terms through
order thirty and the radius contains only the analytic tail and verified
rounding enclosure.

\section{Formal Arb ball certificate}

\subsection{Outward-rounded coefficient extraction}

Arb implements arbitrary-precision midpoint-radius arithmetic with automatic
error tracking \cite{Johansson2017}.  The final audit uses 192-bit real and
complex balls.  Instead of exponentiating large block matrices in ball
arithmetic, amplitude coefficients are extracted by a 64-point Cauchy
discrete Fourier transform on the complex unit circle:
\begin{equation}
\widehat g_n
=\frac{1}{64}
\sum_{k=0}^{63}
g(\zeta_k)\zeta_k^{-n},
\qquad
\zeta_k=e^{2\pi\ii k/64}.
\label{eq:cauchy}
\end{equation}
The discrete transform aliases coefficients $g_{n+64\ell}$.  A path-uniform
Dyson bound encloses this alias.  The implemented enclosure is $10^{-40}$,
and the formal gate verifies that
\[
e^{K_a}\frac{K_a^{64}}{64!}<10^{-40}
\]
for every physical error axis.  Fidelity coefficients are then formed by
ball convolution,
\[
F_n=\sum_{r=0}^n\overline{g_r}g_{n-r},
\]
and the infidelity coefficients follow from $1-F$.

\subsection{Conditional formal theorem}

\begin{theorem}[Certified ordering for the serialized model]
\label{thm:formal}
Consider the twelve controls serialized in the formal certificate with
protocol hash
\begin{center}
\texttt{2c27f222ff6bcb7e87f9073da1b3f151}\\
\texttt{9c08498124e772796d363cab8f5c5338}.
\end{center}
For the finite-dimensional Hamiltonian in Eq.~\eqref{eq:Hamiltonian}, the
common reference target, and the six-axis error law in
Eq.~\eqref{eq:errorlaw}, the 192-bit Arb order-thirty intervals are pairwise
disjoint and certify the complete ordering of all 66 unordered path pairs.
\end{theorem}

\begin{proof}[Computer-assisted proof]
All pulse phases and constants are parsed as serialized decimals or exact
rational multiples of the Arb enclosure of $\pi$.  Arb complex balls enclose
every trigonometric evaluation, matrix exponential, propagation, Cauchy sum,
and coefficient convolution.  The Cauchy alias gate proves that the omitted
Fourier aliases lie inside the declared $10^{-40}$ complex enclosure.  The
analytic even Dyson tail after order thirty is enclosed by $2\times10^{-11}$.
The resulting 12 order-thirty real balls are pairwise disjoint in the order
recorded in the pre-outcome certificate.  Direct interval comparisons certify
all $\binom{12}{2}=66$ pairs.  The certificate is written and hashed before
ordinary held-out outcomes are evaluated.  Finally, every ordinary outcome
is verified to lie inside its corresponding formal ball.
\end{proof}

\begin{remark}[Scope of Theorem~\ref{thm:formal}]
The theorem is conditional on the serialized-decimal finite Hamiltonian
model and the listed candidate controls.  It does not certify model
discrepancy, calibration drift, sampling noise, or PASQAL hardware.
\end{remark}

\begin{theorem}[Exact-root finite-error ordering]
\label{thm:exactroot}
Let $X_k$ be the twelve phase boxes certified in
Theorem~\ref{thm:krawczyk}.  For the declared six-axis mean finite-error
functional, direct 192-bit outward-rounded Arb propagation over the full
boxes $X_k$ produces pairwise disjoint intervals in the prospectively frozen
order.  Hence the locally unique exact matched roots inside the boxes obey
all 66 frozen pairwise orderings.
\end{theorem}

\begin{proof}[Computer-assisted proof]
Each exact root belongs to its certified phase box.  Every trigonometric
evaluation, Hamiltonian matrix exponential, state propagation, and
finite-error infidelity is evaluated over the complete phase enclosure with
outward-rounded Arb/ACB arithmetic.  Direct comparison separates all
$\binom{12}{2}=66$ interval pairs and their orientations equal the frozen
order.  The exact-root ordering protocol and certificate hashes are,
respectively,
\begin{center}
\texttt{6f055d4c33d3af3619123d4e2f87d12d}\\
\texttt{062d232849744037c59a605c86727674},
\end{center}
and
\begin{center}
\texttt{c1aea9572f503221f526b03fbe2b2acb}\\
\texttt{4892d1efb95952eb596eff8b7424ced9}.
\end{center}
\end{proof}

\begin{remark}[Order-thirty mechanism on the root boxes]
Propagating the order-thirty local-jet enclosure over the full Krawczyk boxes
certifies 42 of 66 pairs.  All 42 orientations agree with the frozen order
and none is reversed; interval widening leaves 24 unresolved.  This partial
mechanism certificate does not veto Theorem~\ref{thm:exactroot}, whose direct
finite-error propagation has no Taylor truncation.
\end{remark}

\section{Results}

\begin{table}[t]
\centering
\caption{Representative prospective and formal audits.  ``Pairs'' denotes
certified unordered comparisons.  Candidate sets differ between stages and
between numerical environments; every stage freezes its protocol and
ranking before held-out evaluation.}
\label{tab:results}
\small
\begin{tabular}{p{3.2cm}cccc}
\toprule
Audit & Paths & Spearman $\rho$ & Certified pairs & Result\\
\midrule
Quartic prospective validation
&20&0.998496&--&pass\\
Covariant tensor validation
&12&1.000000&--&pass\\
Order-30 floating certificate
&12&1.000000&66/66&pass\\
Arb formal: quartic only
&12&--&35/66&partial\\
Arb formal: order 30
&12&1.000000&66/66&pass\\
Exact matched roots: Krawczyk
&12&--&12/12&pass\\
Exact-root direct finite error
&12&1.000000&66/66&pass\\
Exact-root order-30 mechanism
&12&--&42/66&partial\\
\bottomrule
\end{tabular}
\end{table}

\begin{table}[t]
\centering
\caption{Numerical integrity diagnostics from representative Colab runs.}
\label{tab:integrity}
\small
\begin{tabular}{lr}
\toprule
Diagnostic & Value\\
\midrule
Fourth-tensor independent components & 15\\
Maximum tensor recovery relative residual & $4.16\times10^{-15}$\\
Direct coordinate-refit relative error & $1.08\times10^{-14}$\\
Invariant-contraction relative error & $1.99\times10^{-15}$\\
Quadratic-coefficient relative spread & $1.33\times10^{-7}$\\
Arb precision & 192 bits\\
Cauchy points & 64\\
Taylor order & 30\\
Formal quartic pair coverage & $35/66=53.03\%$\\
Formal order-30 pair coverage & $66/66=100\%$\\
Strict Krawczyk inclusions & $12/12=100\%$\\
Exact-root direct pair coverage & $66/66=100\%$\\
Exact-root order-30 pair coverage & $42/66=63.64\%$\\
Exact-root order-30 reversed pairs & $0$\\
\bottomrule
\end{tabular}
\end{table}

The sequence of results supports the following hierarchy.
\begin{enumerate}[leftmargin=2.2em]
\item \textbf{Implementation memory:} endpoint and complete first-order
output-state matching do not numerically determine finite-error performance.
\item \textbf{Low-order organization:} the covariant quartic contraction
$\cG_4$ strongly organizes mean performance and formally certifies a
substantial subset of path pairs.
\item \textbf{High-order closure:} close comparisons depend on higher local
derivatives.  The order-thirty local jet plus a bounded tail certifies the
complete ordering for the serialized cohort.
\item \textbf{Exact-root closure:} Krawczyk inclusions certify locally unique
exact matched representatives, and direct phase-box propagation certifies
that all 66 frozen orderings hold at those roots.
\end{enumerate}

\section{Discussion}

\subsection{What is geometrical?}

The geometrical object is not the scalar $\cG_4$ in isolation.  It is the
fourth-order response tensor $A^{(4)}$ together with the physical fourth noise
moment $M^{(4)}$.  Their contraction is independent of a linear reparameter-
ization of error coordinates.  This avoids interpreting a coordinate-specific
axis average as an intrinsic curvature.

The tensor should also not be identified automatically with Riemann
curvature.  It is a quartic susceptibility of the endpoint loss---a component
of the fourth jet of the endpoint map composed with infidelity.  Establishing
a relation to a canonical connection or curvature on a control quotient
would require additional structure.

\subsection{Why the quartic term is useful but incomplete}

First-order matching makes the quadratic loss nearly common.  This promotes
the quartic term to the first visible discriminator.  It explains why
$\cG_4$ often gives almost perfect rankings.  However, if two quartic
coefficients are close, the sixth and higher terms can reverse their order.
The formal intervals expose this distinction directly: 35 pairs are
quartic-certified, while the remaining pairs require higher jet data.

This is preferable to declaring the quartic rule either universally true or
false.  The correct statement is gap-dependent:
\begin{equation}
|\cG_{4,a}-\cG_{4,b}|
>
B_{a}^{(\ge6)}+B_{b}^{(\ge6)}
\quad\Longrightarrow\quad
\text{quartic ordering is certified}.
\label{eq:gap}
\end{equation}

\subsection{Relation to analyticity}

An analytic finite-dimensional propagator is determined by its full local
germ, so the existence of some sufficiently high reconstructing order is not
by itself surprising.  The nontrivial computational content is quantitative:
\begin{itemize}[leftmargin=2em]
\item a finite order, 30, suffices at the full held-out scale;
\item the omitted tail is explicitly bounded;
\item all pairwise intervals are disjoint;
\item a much lower fourth-order object already certifies a substantial
fraction of comparisons.
\end{itemize}
Thus the result is not that ``analytic functions have Taylor series,'' but
that a finite, auditable local resource closes a nontrivial control-ranking
problem.

\subsection{Implications for control selection}

For a noise moment $M$, the low-cost screening problem is
\[
\gamma^\star
\in
\arg\min_{\gamma\in\cF_\eta}
A^{(4)}_\gamma:M^{(4)}.
\]
Pairs that fail the quartic gap test in Eq.~\eqref{eq:gap} can be escalated
adaptively to sixth, eighth, or higher jet order until their intervals
separate.  This suggests a certified branch-and-bound control selector:
compute low-order local tensors for all candidates, refine only ambiguous
pairs, and reserve direct finite-error simulation or hardware evaluation for
model validation rather than exhaustive ranking.

\section{Limitations and next theorem}

The present work has five explicit limitations.
\begin{enumerate}[leftmargin=2.2em]
\item \textbf{Local rather than global fibre theorem.}
The Krawczyk proof certifies a locally unique exact representative inside
each declared transverse box.  It does not establish global uniqueness,
connectedness, or a global eight-dimensional matched manifold.
\item \textbf{Finite-dimensional model.}
Spontaneous emission, Doppler effects, laser phase noise, position
fluctuations, leakage levels, and waveform filtering are absent.
\item \textbf{Two atoms.}
No many-body scaling theorem is claimed.
\item \textbf{Mean axial error law.}
The primary result concerns the declared six-axis mean.  Worst-case ranking
is not supported by the present scalar.
\item \textbf{No QPU evidence.}
The calculations are local and require no PASQAL credentials.  They do not
constitute PASQAL Cloud, FRESNEL, or physical-QPU evidence.
\end{enumerate}

The exact-root finite-radius chain is now closed:
\[
\text{exact matched control}
\;+\;
\text{direct bounded phase-box propagation}
\;\Longrightarrow\;
\text{finite-error ordering}.
\]
The next theorem should incorporate an open-system or experimentally
calibrated uncertainty model and determine whether the certified gaps survive
model discrepancy.  Many-body scaling and independent hardware validation
remain separate problems.

\section{Reproducibility}

The principal standalone scripts are:
\begin{itemize}[leftmargin=2em]
\item \texttt{pasqal\_two\_atom\_G4\_standalone\_colab.py};
\item \texttt{pasqal\_L3\_L4\_standalone\_colab.py};
\item \texttt{pasqal\_L4\_order30\_standalone\_colab.py};
\item \texttt{pasqal\_L4\_formal\_arb\_standalone\_colab.py};
\item \texttt{pasqal\_L4\_exact\_fibre\_krawczyk\_standalone\_colab\_v1\_2.py};
\item \texttt{pasqal\_L4\_exact\_root\_ordering\_standalone\_colab.py}.
\end{itemize}
The final formal protocol hash is
\begin{center}
\texttt{2c27f222ff6bcb7e87f9073da1b3f151}\\
\texttt{9c08498124e772796d363cab8f5c5338}.
\end{center}
The certificate hash is run-specific because the certificate contains a
creation time and the actual accepted path realization.  The formal
standalone script pins \texttt{python-flint==0.8.0}; NumPy and SciPy are used
for path generation, while Arb/ACB is used for the formal interval
certificate.

\section{Conclusion}

We have separated finite-error robustness inside a numerically matched
neutral-atom implementation fibre into a hierarchy of local information.
The common endpoint and complete first-order state response leave residual
path dependence.  A coordinate-covariant fourth-order response tensor
provides a strong and partially certifying low-order description.  Close
comparisons require higher derivatives.  For a frozen twelve-path
serialized model, the zero-error jet through order thirty, combined with
rigorous Cauchy-alias and Dyson-tail enclosures, yields an Arb
outward-rounded certificate of all 66 finite-error pairwise orderings.
Krawczyk inclusions then certify locally unique exact matched roots in all
twelve declared transverse boxes, and direct Arb propagation of those boxes
certifies the same 66/66 order at the exact roots.

The strongest justified conclusion is therefore conditional but precise:
finite-error ordering in this model is recoverable from a finite local jet,
and the ordering persists at locally unique exact matched representatives
under direct formal phase-box propagation.  Global fibre structure and
extensions to open-system, many-atom, model-discrepancy, and hardware data
remain open.

\begin{thebibliography}{99}

\bibitem{Silverio2022}
H.~Silv\'erio, S.~Grijalva, C.~Dalyac, L.~Leclerc,
P.-J.~Karalekas, N.~Shammah, M.~Beji, L.-P.~Henry, and L.~Henriet,
``Pulser: An open-source package for the design of pulse sequences in
programmable neutral-atom arrays,''
\emph{Quantum} \textbf{6}, 629 (2022).
\href{https://doi.org/10.22331/q-2022-01-24-629}
{doi:10.22331/q-2022-01-24-629}.

\bibitem{Henriet2020}
L.~Henriet, L.~Beguin, A.~Signoles, T.~Lahaye, A.~Browaeys,
G.-O.~Reymond, and C.~Jurczak,
``Quantum computing with neutral atoms,''
\emph{Quantum} \textbf{4}, 327 (2020).
\href{https://doi.org/10.22331/q-2020-09-21-327}
{doi:10.22331/q-2020-09-21-327}.

\bibitem{Saffman2016}
M.~Saffman,
``Quantum computing with atomic qubits and Rydberg interactions:
progress and challenges,''
\emph{J. Phys. B: At. Mol. Opt. Phys.} \textbf{49}, 202001 (2016).
\href{https://doi.org/10.1088/0953-4075/49/20/202001}
{doi:10.1088/0953-4075/49/20/202001}.

\bibitem{Wimperis1994}
S.~Wimperis,
``Broadband, narrowband, and passband composite pulses for use in
advanced NMR experiments,''
\emph{J. Magn. Reson. A} \textbf{109}, 221--231 (1994).
\href{https://doi.org/10.1006/jmra.1994.1159}
{doi:10.1006/jmra.1994.1159}.

\bibitem{Kukita2022}
S.~Kukita, H.~Kiya, and Y.~Kondo,
``General off-resonance-error-robust symmetric composite pulses with
three elementary operations,''
\emph{Phys. Rev. A} \textbf{106}, 042613 (2022).
\href{https://doi.org/10.1103/PhysRevA.106.042613}
{doi:10.1103/PhysRevA.106.042613}.

\bibitem{Barnes2015}
E.~Barnes, X.~Wang, and S.~Das Sarma,
``Robust quantum control using smooth pulses and topological winding,''
\emph{Sci. Rep.} \textbf{5}, 12685 (2015).
\href{https://doi.org/10.1038/srep12685}
{doi:10.1038/srep12685}.

\bibitem{Johansson2017}
F.~Johansson,
``Arb: Efficient arbitrary-precision midpoint-radius interval
arithmetic,''
\emph{IEEE Trans. Comput.} \textbf{66}, 1281--1292 (2017).
\href{https://doi.org/10.1109/TC.2017.269063}
{doi:10.1109/TC.2017.269063}.

\bibitem{Dyson1949}
F.~J.~Dyson,
``The radiation theories of Tomonaga, Schwinger, and Feynman,''
\emph{Phys. Rev.} \textbf{75}, 486--502 (1949).
\href{https://doi.org/10.1103/PhysRev.75.486}
{doi:10.1103/PhysRev.75.486}.

\bibitem{Tucker2011}
W.~Tucker,
\emph{Validated Numerics: A Short Introduction to Rigorous Computations},
Princeton University Press, Princeton (2011).

\bibitem{Moore2009}
R.~E.~Moore, R.~B.~Kearfott, and M.~J.~Cloud,
\emph{Introduction to Interval Analysis},
SIAM, Philadelphia (2009).

\end{thebibliography}

\end{document}
